\begin{table}[htbp]
\centering
\caption{HCPCS Code Classification: T-Codes (HCBS) by In-Home Status}
\label{tab:hcpcs}
\small
\begin{tabular}{llrrr}
\hline\hline
Code & Description & Spending (\$B) & Claims (M) & In-Home \\
\hline
T1000 & Other T-code & 7.0 & 16.7 & No \\
T1003 & Other T-code & 2.8 & 12.4 & No \\
T1015 & FQHC visit & 49.0 & 320.8 & No \\
T1016 & Other T-code & 5.0 & 94.2 & No \\
T1017 & Other T-code & 8.4 & 74.7 & No \\
T1018 & Other T-code & 2.4 & 32.5 & No \\
T1019 & Personal care, 15 min & 121.7 & 1086.3 & Yes \\
T1020 & Personal care, per diem & 8.2 & 33.5 & Yes \\
T1040 & Other T-code & 5.7 & 29.9 & No \\
T2003 & Non-emergency transport & 2.4 & 104.5 & No \\
T2013 & Other T-code & 2.0 & 12.2 & No \\
T2016 & Habilitation residential, per diem & 34.0 & 65.5 & Yes \\
T2017 & Habilitation residential, 15 min & 2.8 & 13.7 & Yes \\
T2021 & Other T-code & 8.6 & 55.2 & No \\
T2022 & Case mgmt, per month & 3.7 & 17.6 & Yes \\
T2023 & Other T-code & 5.1 & 12.3 & No \\
T2025 & Waiver services NOS & 3.6 & 20.3 & Yes \\
T2031 & Other T-code & 4.5 & 13.3 & No \\
T2033 & Supported employment & 7.4 & 16.5 & Yes \\
T2046 & Other T-code & 4.2 & 14.1 & No \\
\hline\hline
\end{tabular}
\begin{minipage}{0.85\textwidth}
\vspace{4pt}
\footnotesize
\textit{Notes:} Top 20 T-prefix HCPCS codes by cumulative spending (2018--2024) from T-MSIS. ``In-Home'' = Yes indicates codes classified as genuinely in-person home-based care in the clean HCBS subset. Codes classified as ``No'' include clinic-based visits (T1015), transportation (T2003, T2005, T2007), and facility-based screening/assessment (T1023, T1024). The clean HCBS subset is the primary classification used in all main analyses.
\end{minipage}
\end{table}
